\section{Why Full FP4 Is Difficult}

\subsection{P lies on the critical path}

For one query row, exact attention is
\begin{equation}
  O_i=\frac{\sum_j e^{z_{ij}}V_j}{\sum_j e^{z_{ij}}},
  \qquad z_{ij}=Q_iK_j^\mathsf{T}/\sqrt d.
  \label{eq:exact-attention}
\end{equation}
FlashAttention visits the keys in tiles. After a tile with maximum
$\widetilde m_i$ arrives, it updates a running maximum $m_i$, denominator
$\ell_i$, and output $o_i$:
\begin{align}
  m_i' &= \max(m_i,\widetilde m_i), \\
  \ell_i' &= e^{m_i-m_i'}\ell_i+
    \sum_{j\in B}e^{z_{ij}-m_i'}, \\
  o_i' &= e^{m_i-m_i'}o_i+
    \sum_{j\in B}e^{z_{ij}-m_i'}V_j.
  \label{eq:online-recurrence}
\end{align}
This avoids storing full score and probability matrices in HBM, but it also
means P is not an input. For full FP4 attention, the first useful PV command
waits for
\begin{equation}
  S\rightarrow\text{maximum}\rightarrow\text{scale}
  \rightarrow\text{probability}\rightarrow\text{E2M1 pack}
  \rightarrow\text{publish}.
  \label{eq:serial-p-path}
\end{equation}
Work removed after publication may not change latency. Work removed before
the first legal P tile can.

On data-center Blackwell GPUs, asynchronous \tcgen{} operations accumulate
FP32 score and output tiles in TMEM. TMA supplies shared-memory operands, and
specialized warps issue matrix operations, compute softmax, and correct the
online output \cite{nvidia2026ptx,zadouri2026flashattention4}. Matrix
throughput grew faster than score-load, exponential, and synchronization
throughput. FP4 accelerates both matrix products while adding P conversion
between them, so this imbalance is stronger than in BF16 FA4.

\subsection{NVFP4 and MXFP4 solve different numerical problems}

Both formats use signed E2M1 payload values
\begin{equation}
  \mathcal F_{\mathrm{E2M1}}=
  \left\{0,\tfrac12,1,\tfrac32,2,3,4,6\right\}
\end{equation}
plus sign. Their scales differ:
\begin{table}[htbp]
\centering
\caption{Block-scaled FP4 formats used in this work.}
\label{tab:fp4-formats}
\small
\begin{tabularx}{\textwidth}{lcccX}
\toprule
Format & Block & Scale & Main strength & Retained use \\
\midrule
NVFP4 & 16 & E4M3, optional tensor scale & fine local placement & signed Q and K \\
MXFP4 & 32 & E8M0 power of two & wide exponent range & nonnegative P and V \\
\bottomrule
\end{tabularx}
\end{table}

For an NVFP4 P block with maximum $a_B$ and global encode factor $G$, the
decoded block scale is approximately
\begin{equation}
  \widehat s_{\mathrm{NV}}=
  \frac{Q_{\mathrm{E4M3}}(G a_B/6)}{G}.
  \label{eq:nv-scale}
\end{equation}
If the E4M3 scale rounds to zero, the entire block disappears. A row-level
stabilizer fixes the range, but its scale and inverse correction add state to
the online path. MXFP4 instead selects a nearby power of two,
\begin{equation}
  \widehat s_{\mathrm{MX}}\in
  \left\{2^{\lfloor\log_2(a_B/6)\rfloor},
          2^{\lceil\log_2(a_B/6)\rceil}\right\},
  \label{eq:mx-scale}
\end{equation}
which folds naturally into a base-two score transform and matches one N32
producer fragment. The choice is a latency--range trade-off, not a claim that
E8M0 is intrinsically more accurate than E4M3.

Table~\ref{tab:p-range} isolates that distinction from kernel scheduling. It
forms exact Gaussian softmax probabilities, quantizes matched N32 blocks,
renormalizes the represented P, and multiplies by the same FP32 V. Stabilized
NVFP4 is the high-fidelity control; unscaled NVFP4 exposes underflow.

\begin{table}[htbp]
\centering
\caption{Probability-format range at D128. ``Zero scales'' is the fraction
of blocks whose scale encodes as zero; ``lost mass'' is exact probability
mass mapped to zero. This is a numerical diagnostic, not a kernel timing.}
\label{tab:p-range}
\scriptsize
\begin{adjustbox}{max width=\textwidth}
\begin{tabular}{llrrrrrr}
\toprule
$S$ & Format & Zero scales & Zero payload & Lost mass & $P$ rel.-$L_2$
& $PV$ cosine & $PV$ rel.-$L_2$ \\
\midrule
1024 & NVFP4 G{=}1 & 0.706299 & 0.831829 & 0.683773 & 1.715156 & 0.700819 & 1.685442 \\
1024 & NVFP4 G{=}448 & 0.000000 & 0.155418 & 0.029314 & 0.114408 & 0.993854 & 0.113687 \\
1024 & MXFP4 & 0.000000 & 0.188683 & 0.040178 & 0.146380 & 0.989846 & 0.144944 \\
4096 & NVFP4 G{=}1 & 0.996063 & 0.999438 & 0.994332 & 13.815642 & 0.217328 & 13.876892 \\
4096 & NVFP4 G{=}448 & 0.000000 & 0.156454 & 0.029760 & 0.114091 & 0.993747 & 0.114582 \\
4096 & MXFP4 & 0.000000 & 0.189755 & 0.040494 & 0.147209 & 0.989338 & 0.148716 \\
8192 & NVFP4 G{=}1 & 0.999634 & 0.999977 & 0.999284 & 12.467307 & 0.096774 & 12.389232 \\
8192 & NVFP4 G{=}448 & 0.000000 & 0.160613 & 0.030691 & 0.114725 & 0.993841 & 0.113660 \\
8192 & MXFP4 & 0.000000 & 0.194189 & 0.041630 & 0.147318 & 0.989660 & 0.146237 \\
\bottomrule
%
\end{tabular}
\end{adjustbox}
\end{table}

\subsection{Relation to prior low-precision attention}

HAO's FP4 FA4 repository is the structural starting point of this work
\cite{hao2026fp4flashattention}. It supports block-scaled FP4 QK with BF16 or
FP8 PV and a stabilized full-NVFP4 path. Its implementation and the Attn-QAT
study identify online P quantization and scale movement as the central cost
\cite{zhang2026attnqat}.

SageAttention3 contributes Q/K smoothing, two-level P scaling, and reuse of
the block maximum during online softmax \cite{zhang2025sageattention3}. Its
RTX 5090 target is SM120, which does not expose the data-center \tcgen{} TMEM
family used by SM100 and SM103 \cite{nvidia2026ptx}. We therefore borrow its
numerical lessons but use HAO and FA4 as the scheduling references. Our focus
is forward kernel performance; we do not claim the training recovery shown by
Attn-QAT.
